\section{Correctness Properties}
\label{sec:correctness}

This section establishes three correctness properties for the Recursive Spawning
mechanism: drift prevention, chain integrity, and termination. Each property is
derived formally from the three-layer integration specified in
Section~\ref{sec:integration}. The properties are not conditional on the honesty
or correctness of any machine actor — they are enforced by the structural
properties of the three-layer architecture and hold regardless of what any agent
or node attempts to do at runtime.

% ---------------------------------------------------------------
\subsection{Drift Prevention}
\label{sec:drift-prevention}

\textit{Drift} in a recursive spawning context refers to the phenomenon by which
a child node's behavior diverges from the principal's intent — either by
expanding its own authority beyond what $R(\text{parent})$ permits, by modifying
the accountability chain, or by replacing the human anchor with a machine actor
as its escalation terminus. Drift is heritable: once an agent enters the drifted
state, all of its children inherit the condition absent an explicit structural
intervention at spawn time. A single undetected drift event at depth $d$ with
branching factor $b$ produces up to $b^d$ drifted agents in the field. The BX3
Framework's three-layer integration is designed to prevent each possible drift
pathway before it can propagate.

\begin{theorem}[Drift Prevention]
\label{th:drift-prevention}
Let $T = (V, E)$ be a BX3 spawn tree satisfying all invariants of
Section~\ref{sec:integration} (parent-pointer immutability, hierarchy finiteness,
delegation scope inheritance, and forensic completeness). Then for any node
$v \in V$, the accountability chain from $v$ to its human anchor $h(v)$ cannot
be redirected, truncated, or absorbed by any machine actor operating within the
BX3 architecture.
\end{theorem}

\begin{proof}[Proof sketch]
The proof proceeds by examining the four distinct pathways by which drift could
occur, and establishing that each is structurally precluded by the three-layer
integration.

\textbf{Pathway 1: Parent-pointer modification.} Suppose a machine actor
attempts to redirect $\alpha(v)$ from its legitimate parent $p$ to a different
node $p'$. By Invariant~\ref{inv:parent-pointer}, $\alpha(v)$ is a read-only
property of the \fact{} layer worksheet, enforced by the pre-commit hook. Any
write attempt is intercepted and rejected before the state advances. The
rejection event is recorded in the Forensic Ledger as
\texttt{parent-pointer-write-denied}$(v, p, p')$, and the Bailout Protocol is
triggered. The drift attempt is both prevented and recorded. No machine actor
has the authority to override this check.

\textbf{Pathway 2: Human anchor replacement.} Suppose a machine actor attempts
to replace $h(v)$ with a machine actor $m$ as the escalation terminus. By
Invariant~\ref{inv:human-anchor}, $h(v)$ is established at node provisioning and
is never modified during the node's operational lifetime. The human anchor is a
structural property of the \purpose{} layer, not a configurable parameter. Any
attempt to modify it is rejected by the \fact{} layer's pre-commit hook, and the
Bailout Protocol is triggered. The drift attempt is prevented.

\textbf{Pathway 3: Authority expansion through spawning.} Suppose a machine
actor at node $p$ attempts to spawn a child $c$ with $R(c) \not\subset R(p)$,
thereby expanding the delegated authority beyond what the principal sanctioned.
By Invariant~\ref{inv:scope-inheritance}, the \bounds{} layer evaluates
$R(c) \subset R(p)$ before every spawn event. If the check fails, the
$e_{\text{scope-violation}}$ event is emitted, the spawning state machine
transitions to \texttt{SPAWN\_LOCKED}, and the Bailout Protocol is triggered.
The attempted authority expansion is prevented and recorded.

\textbf{Pathway 4: Chain truncation via non-termination.} Suppose a machine
actor attempts to prevent the escalation chain from reaching $h(v)$ by causing
the chain to loop or diverge. By Invariant~\ref{inv:hierarchy-finiteness}, the
parent-pointer chain from any node has maximum length $D_{\max}$, and by
Invariant~\ref{inv:spawn-accountability} it terminates at a human principal. The
chain cannot loop because $\alpha(v)$ is immutable and acyclic — a node cannot
be its own ancestor. The chain cannot terminate at a machine actor because
Invariant~\ref{inv:human-anchor} forbids the human anchor from being replaced.
The drift attempt is structurally impossible.

Since all four drift pathways are precluded, the accountability chain from any
node $v$ to its human anchor $h(v)$ is non-divertable. $\square$
\end{proof}

The drift prevention theorem depends critically on the conjunction of all three
layers operating in concert. No single constraint is sufficient: parent-pointer
immutability alone does not prevent authority expansion; delegation scope
inheritance alone does not prevent chain redirection; hierarchy finiteness
alone does not prevent the human anchor from being replaced. The Forensic Ledger
plays an essential evidentiary role — even though the pre-commit hooks prevent
drift at runtime, the ledger ensures that any attempted drift is permanently
recorded with cryptographic chain integrity. The combination of prevention and
attribution is what makes the upstream accountability guarantee operational in
high-stakes deployments.

% ---------------------------------------------------------------
\subsection{Chain Integrity}
\label{sec:chain-integrity}

\textit{Chain integrity} refers to the guarantee that the parent-pointer chain
from any node to its human anchor is complete, tamper-evident, and
reconstructable. The chain is the escalation path used by the Bailout Protocol;
if the chain is incomplete or tampered with, the escalation path is unreliable
and the liveness guarantee is compromised.

\begin{property}[Chain Integrity]
\label{prop:chain-integrity}
For any node $v$ in a BX3 spawn tree, the accountability chain
$C(v) = \langle v, \alpha(v), \alpha^{(2)}(v), \ldots, \alpha^{(k)}(v) \rangle$
where $\alpha^{(k)}(v) = h(v)$ satisfies:
\begin{enumerate}
\item $k \leq D_{\max}$ (finite length);
\item $\alpha^{(i)}(v)$ is a machine node for all $0 \leq i < k$;
\item $\alpha^{(k)}(v)$ is a human principal;
\item For all $0 \leq i < j \leq k$, $\alpha^{(i)}(v) \neq \alpha^{(j)}(v)$ (acyclic);
\item Every entry in $C(v)$ appears as a \texttt{child-provisioned} event in the
Forensic Ledger with matching $\alpha(\cdot)$ and $h(\cdot)$ values.
\end{enumerate}
\end{property}

Items (1)–(3) follow from the hierarchy finiteness constraint, the delegation
scope inheritance constraint, and the human anchor invariant. Item (4) follows
from the parent-pointer immutability invariant: since $\alpha(v)$ is established
once at node creation and never modified, and since $\alpha(v) \neq v$ for all
$v$, the parent-pointer relation defines a strict partial order with no cycles.
Item (5) follows from the Forensic Ledger completeness invariant: every spawn
event is recorded before subsequent actions are permitted, and every
\texttt{child-provisioned} entry encodes both $\alpha(c) = p$ and $h(c) = h(p)$.

The tamper-evidence property follows from the Forensic Ledger's cryptographic
chaining. Each entry carries the SHA-256 hash of the preceding entry. Any
modification to an entry — whether to the node identifier, the event type, the
context, or the parent pointer — breaks the chain at the point of modification.
The break is detectable by recomputing the chain of hashes from the genesis
entry to the entry in question and comparing against the stored hash. An auditor
without access to the ledger's contents can detect tampering with high
confidence.

The reconstructability property follows from the combination of (5) and the
cryptographic chaining. Given any node $v$, an auditor can reconstruct $C(v)$
by querying the Forensic Ledger for all \texttt{child-provisioned} entries
involving $v$, extracting the $\alpha(v)$ field from each entry, and traversing
the chain until reaching a node that does not appear as a child in any
\texttt{child-provisioned} entry — that node is $h(v)$. The reconstruction is
deterministic and reproducible; any two auditors who independently reconstruct
$C(v)$ will obtain identical results.

% ---------------------------------------------------------------
\subsection{Termination}
\label{sec:termination}

\textit{Termination} is the guarantee that the escalation chain from any node
always reaches a human principal in finite time, and that the Bailout Protocol's
escalation algorithm terminates. This is the structural foundation of the BX3
Framework's Liveness guarantee.

\begin{property}[Finite Termination]
\label{prop:finite-termination}
For any node $v$ in a BX3 spawn tree, the escalation chain $C(v)$ terminates at
a human principal $h(v)$ in at most $D_{\max}$ steps.
\end{property}

\begin{proof}[Proof sketch]
By the hierarchy finiteness constraint, every node $v$ has depth
$d(v) \leq D_{\max}$, where $d(\text{root}) = 0$ and $d(v) = d(\alpha(v)) + 1$
for all $v \neq \text{root}$. The root node is provisioned directly by the
human principal and therefore satisfies $h(\text{root}) = \text{principal}$.
For any node $v$, the chain $v, \alpha(v), \alpha^{(2)}(v), \ldots,
\alpha^{(d(v))}(v)$ exists and has length $d(v) \leq D_{\max}$. By
Invariant~\ref{inv:human-anchor}, $\alpha^{(d(v))}(v)$ is a human principal.
Therefore $C(v)$ terminates at a human principal in at most $D_{\max}$ steps.
$\square$
\end{proof}

\begin{property}[Deterministic Escalation Path]
\label{prop:deterministic-escalation}
For any node $v$ and any exception state $e$ at $v$, the Bailout Protocol's
escalation path from $v$ to $h(v)$ is uniquely determined and is independent of
the content of $e$ or the state of the \bounds{} layer at any node on the path.
\end{property}

Property~\ref{prop:deterministic-escalation} follows from the immutability of
the parent-pointer chain. The escalation algorithm traverses
$\alpha(v), \alpha^{(2)}(v), \ldots, \alpha^{(k)}(v)$ where $k \leq D_{\max}$.
Each $\alpha^{(i)}(v)$ is immutable and deterministic — established once at node
creation and never modified. The path does not depend on the exception type, the
exception severity, or any runtime state. It is a structural property of the
spawn tree, not an operational decision made by any machine actor. This means
that the escalation path is predictable and auditable before any exception
occurs: the human anchor can determine the complete escalation path for any
node by reading the Forensic Ledger entries for all \texttt{child-provisioned}
events along the chain.

Together, drift prevention, chain integrity, and finite termination constitute
the full correctness specification for the Recursive Spawning mechanism. The
spawn tree is a structurally sound representation of the delegation chain: every
edge represents a genuine, auditable, bounded delegation of authority from a
principal (human or machine) to a child node, and the chain from any leaf node
to its human anchor is fully recorded, finite in depth, tamper-evident, and
non-circumventable.